\chapter{结论与展望}

\section{结论}
本文研究不确定环境下机器人的安全强化学习问题，按"单体安全—群体协同—系统验证"三条线索展开：第3章解决单机器人执行层的硬安全保障，第4章把安全约束推广到多机器人编队的几何关系，第5章将两者集成到统一的容器化平台上进行验证。

第3章提出约束流形策略学习方法（CMPL）。其核心是把策略输出投影到安全可达空间的切空间，使动作在执行前即被修正到满足硬约束，而非依赖奖励惩罚的间接引导；状态估计则用一个由卷积神经网络学习观测噪声协方差的扩展卡尔曼滤波器，为安全过滤器提供更可靠的输入。在 Safety-Gymnasium 的目标到达任务中，CMPL 将平均安全代价从 PPO 基线的 50.58 降到 4.91（降幅 90.3\%）；在 Webots 的 E-puck 零样本部署中，碰撞率由 40\% 降至零，位置误差降低 27.5\%。

第4章针对多机器人编队中任务目标与安全约束相互耦合的问题，提出几何约束策略学习方法（GCPL）。该方法用 MAPPO 学习任务级控制意图，再由几何过滤层在黎曼流形上统一处理编队拓扑、机间安全距离与避障，并对动作做解析修正，把任务驱动与结构约束解耦到两层。在线形编队任务上，GCPL 相对 MAPPO 基线将安全代价从 105.78 降到 18.53，编队完成度由 10\% 提升到 71\%，成功率由 2\% 提升到 66\%；在楔形、圆形等更复杂拓扑上同样保持领先。

第5章面向方法验证与工程落地，设计并实现了一套容器化仿真验证平台：用分层架构和 Docker 镜像固化运行依赖，并提供统一的评估管线支持多算法在相同条件下的公平对比。平台把第3、4章方法按 MAPPO→GCPL→CMPL 的顺序串联，先在纯 Python 的 MVP 上验证控制链路的数学正确性，再迁移到 Webots 做四档消融。四档实验全程零碰撞，完整分层配置在编队一致性上最优（编队误差 0.082~m），到达率则随各模块权重组合在安全裕度与目标驱动之间权衡——这说明安全机制能在策略尚未收敛时提供保护，但任务完成度仍取决于策略本身的优化程度。

三部分工作由"在动作执行前于约束流形上显式修正安全约束"这一共同思想衔接：CMPL 给出单体硬安全，GCPL 把它推广到群体几何协调，平台则让二者在统一环境下端到端运行并可复现地对比。

\section{展望}

尽管本文在多机器人安全强化学习方面取得了一定进展，但仍存在若干值得进一步研究的问题，未来可从以下几个方向开展深入探索：

（1）本文方法主要依赖离线可达性分析与先验环境信息，在动态障碍物或未知环境中仍存在一定局限。未来可结合在线学习或自适应可达性分析方法，实现安全区域的动态更新，从而提升系统在强不确定环境中的适应能力。

（2）随着机器人数量的增加，多智能体系统在通信负载、计算复杂度和策略协同方面会面临更高要求。未来可结合图神经网络、分布式学习等方法，进一步提升算法在大规模系统中的扩展能力与计算效率。

（3）本文实验主要基于仿真环境完成，后续工作仍需在实际机器人平台上开展进一步实验，重点考察传感器误差、执行器延迟、通信不稳定和场景扰动等因素对方法性能的影响。

（4）本文提出的 GCPL 方法在常规编队导航任务中能够降低碰撞风险并提升队形稳定性，但在密集编队、狭窄通道等复杂场景中，固定权重或度量矩阵可能导致局部极小值、路径震荡等问题。后续可引入动态权重调整、任务优先级和冲突重规划机制，使编队能够根据环境变化临时调整形态，并在通过复杂区域后恢复目标队形。

（5）本文的数据驱动 EKF 主要基于速度相关的高斯噪声模型进行训练，能够刻画运动速度增大导致测量误差上升的情况。但真实环境中的噪声可能更加复杂。后续可进一步考虑非高斯噪声、突发异常和分布外场景，并结合异常检测、在线校准或不确定性估计，提高状态估计模块在真实环境中的稳定性。